\documentclass[12pt]{article}
\usepackage[obeyspaces]{url}
\usepackage{amsmath,amssymb}
\usepackage{fullpage}
\usepackage{hyperref}
\usepackage{mathtools}
\usepackage{seqsplit}
\usepackage{tikz}
\usepackage{MnSymbol,wasysym}
\usepackage{pdflscape}
\usepackage{lscape}
\usepackage{newpxtext}
\usepackage{booktabs}

\setlength{\parindent}{0pt}
\setlength{\parskip}{2ex}
\thispagestyle{empty}

\hypersetup{
    colorlinks=true,
    linkcolor=blue,
    filecolor=magenta,      
    urlcolor=cyan,
    pdftitle={Overleaf Example},
    pdfpagemode=FullScreen,
    }

\newcommand{\N}{\ensuremath{\mathbb{N}}}
\newcommand{\Q}{\ensuremath{\mathbb{Q}}}
\newcommand{\R}{\ensuremath{\mathbb{R}}}
\newcommand{\Z}{\ensuremath{\mathbb{Z}}}

\begin{document}

\includegraphics{CSUSig}

\begin{center}
    {\large\textbf{MATH 101 - Section 801 - Final Poster Project}}
\end{center}
\textbf{\underline{Instructions:}} Instead of a final exam, you will prepare a poster about a topic related to math in the social sciences. You will participate in our virtual Math 101 Poster Session during Final Exam week (\textbf{05/12-05/17}) on Canvas, where you will interact with your peers on a Canvas discussion board. Key dates are as follows:

\begin{itemize}
    \item \textbf{Friday 05/02}, 11:59pm: Submit your project topic on Canvas.
    \item \textbf{Friday 05/09}, 11:59pm: Upload your poster as a \textbf{.pdf} to the Canvas assignment.
    \item \textbf{Friday 05/12}: I will make all posters available for viewing on a Canvas discussion board.
    \item \textbf{Wednesday 05/14}, 6pm: Reply to 3 different posters describing something you found interesting, what you liked about the poster, and one question you still have.
\end{itemize}

\textbf{\underline{Format:}} Final poster dimensions can be up to a maximum of 48'' wide x 36'' tall landscape format or 36'' wide x 48'' tall vertical format (although the posters will not be printed, these dimensions impact the scale of your materials.) Your final poster should be saved as a \textbf{.pdf} file. Your poster should include your name, poster title, and references.

\textbf{\underline{Topic:}}
You can take any topic that we have covered and find \textbf{real-world} examples. You can do a poster on topics in the book we \textbf{haven't} covered. Finally, you can find math in something that \textbf{interests you} (music, sports, art, biology, weather, finances, crime, AI, space travel, etc.). Google ``math in...'' to see what comes up. Questions to think about:
\begin{itemize}
    \item What problems were people attempting to solve?
    \item What background or math do we need to know to understand this topic?
    \item How did society (people, current events, technology, etc.) influence this topic?
\end{itemize}

You will submit a project proposal to me on \textbf{Friday 05/02} with your chosen topic, and I will approve it or offer suggestions. If you tell me your interests, I am happy to offer topic ideas!

\textbf{\underline{Poster Advice}}
The principles for a poster are similar to those of a presentation, meaning you should prioritize pictures and short descriptions. Unlike a presentation, however, a poster needs to be understood without you there, so some of your descriptions can be a bit longer than what you would use for a presentation.

Here is some advice taken from Anya Michaelsen's excellent \href{https://www.ams.org/journals/notices/202002/rnoti-p189.pdf}{article} on the topic:
\begin{itemize}
    \item \textbf{Structure your Poster Effectively:} Organize your poster into clear and defined sections, using boxes and titles. Follow a logical flow, typically from top to bottom and left to right. Consider sections like Introduction, Background, Results, and Discussion.

    \item \textbf{Use White Space:} Avoid overwhelming your poster with dense text. Keep text concise and use bullet points or lists. Utilize headings, subheadings, and space between elements for clarity.

    \item \textbf{Incorporate Visual Elements:} Use color, bold, italics, and other formatting to highlight key points. Include figures, diagrams, or pictures to break up text and enhance understanding. Consider color coding to differentiate important concepts.

    \item \textbf{Tailor Content to Your Audience:} Determine your intended audience's background and knowledge level. Adjust the level of detail in your background section accordingly. Ensure your poster can stand alone if necessary or have an accompanying one-minute summary.
\end{itemize}
\newpage
The following can help you get started on your poster:
\begin{table*}[h!]
    \centering
    \begin{tabular}{p{0.45\linewidth}p{0.45\linewidth}}\toprule
        \textbf{PowerPoint} & \textbf{Goggle Slides}\\ \midrule
        \begin{enumerate}
            \item In PowerPoint, begin by creating a 1-slide presentation. Choose a blank slide. (The entire poster must be contained in only one slide.)
            \item Adjust the slide size to make the poster dimensions to a maximum of $36"\times48"$ or $48"\times36"$ or a smaller size if desired.
            \item Save your poster as a .pptx file while you're working on it. (When you are finished and ready to submit, save a copy as a .pdf)
        \end{enumerate}&
        \begin{enumerate}
            \item Go to slides.google.com in your browser.
            \item Click the big ``+'' symbol to start a new slideshow.
            \item In ``File/Page Setup'' choose ``Custom'' and set the size of your slide to $36"\times48"$ or $48"\times36"$ or to a smaller size if desired.
            \item In ``Slide/Apply Layout'' choose ``Blank''.
            \item When you're finished and ready to submit, choose ``File/Download'' and choose ``PDF document''.
        \end{enumerate}\\
   \bottomrule
    \end{tabular}
    \end{table*}

\underline{\textbf{Templates:}} You don't have to design your poster from scratch. You can, and probably should, use a template to get started. Spend time thinking about how you want your poster to look before choosing a template. Below I linked some resources for templates, but there are many other places on the internet to get templates.
\begin{itemize}
    \item University of Minnesota: \url{https://libguides.umn.edu/posters/googleslides}
    \item Western Carolina University: \url{https://researchguides.wcu.edu/researchposter/layout}
    \item The Craft of Scientific Presentations: \url{https://www.craftofscientificposters.com/templates.html}
    \item PosterPresentations.com: \url{https://www.posterpresentations.com/how-to-add-our-poster-templates-to-google-slides.html}
\end{itemize}

\underline{\textbf{Examples:}} You will find examples on Canvas. Take note of things you like and don't like.

\textbf{\underline{Grading Rubric:}} With a total of 18 points, please see next page.
\begin{table*}[h!]
    \centering
    \begin{tabular}{p{0.4\linewidth}p{0.4\linewidth}p{0.05\linewidth}}\toprule
        \textbf{Criteria} & \textbf{Marks} & \textbf{Pts.}\\ \midrule
        {\textbf{Clarity of poster}:\par Easily legible, clear flow of ideas, has headings, images, title, etc.}& 3 (meets expectations)\par 2 (near expectations)\par 1 (below expectations)\par 0 (missing) & 3\\\midrule
        \textbf{Math}:\par Math involved is correct,  explained and meaningful to the topic& 3 (meets expectations)\par 2 (near expectations)\par 1 (below expectations)\par 0 (missing) & 3\\\midrule
        \textbf{Background}:\par Math or other content needed to understand the rest of the poster is explained & 3 (meets expectations)\par 2 (near expectations)\par 1 (below expectations)\par 0 (missing) & 3\\\midrule
        \textbf{References}:\par References are included on poster, can be cited in any format you wish. &3 (meets expectations)\par 2 (near expectations)\par 1 (below expectations)\par 0 (missing) & 3\\\midrule
        \textbf{Peer poster responses}:\par Write responses on $3$ posters you saw at  virtual poster session.& 3 (meets expectations)\par 2 (near expectations)\par 1 (below expectations)\par 0 (missing) & 3\\\midrule
        \textbf{Timely submission of proposal}:\par Submitted the project proposal by the due date with a small explanation on the main idea& 2 (meets expectations)\par 1 (below expectations)\par 0 (Missing)& 2\\\midrule
        \textbf{Timely submission of poster}:\par Submitted the poster by the due date and in the correct format.& 1 (did)\par 0 (didn't)& 1\\
   \bottomrule
    \end{tabular}
    \end{table*}
\end{document}